\section{Technical proofs for nonexceptional endpoints}
\label{app:nonexceptional-proofs}

\subsection{Proofs supporting the one-dimensional reduction}
\label{app:local-reduction}

\begin{proof}[Proof of \Cref{lem:boundary-convergence}]
Put
\[
  q_n=\pc(r_n),
  \qquad
  p_0=\pc(r_0).
\]
By continuity of $\pc$ at $r_0$, we have $q_n\to p_0$, and
$|p_n-q_n|\to0$ gives $p_n\to p_0$ as well.
Choose a compact interval $P\subset(0,1)$ with $p_0$ in its interior.
After discarding finitely many terms, both $p_n$ and $q_n$ lie in $P$.
The role of $P$ is only to keep the logarithms in $I_p$ uniformly away from the
singular endpoints $p=0,1$.

Since $(\widetilde{\W}_0,\delta_\cut)$ is compact, every subsequence of
$(W_n)$ has a further cut-convergent subsequence; so it suffices to show that
every cut-convergent subsequence has limit $r_0$.
Passing to such a subsequence and relabeling it, we assume that $W_n\to W_*$ in cut distance.
We will prove that $W_*\equiv r_0$.
First note that the homomorphism densities
are continuous in the cut metric, so
\begin{equation}
  \label{eq:boundary-limit-feasibility}
  t(H,W_*)\ge r_0^m.
\end{equation}

We next compare the relative-entropy functionals with nearby values of the
parameter $p$.  For $p,q\in P$, the identity
\[
  I_p(W)
  =
  -2s(W)-\log(1-p)+e(W)\log\frac{1-p}{p}
\]
gives
\[
  I_p(W)-I_q(W)
  =
  \log\frac{1-q}{1-p}
  +e(W)\log\frac{(1-p)q}{p(1-q)}.
\]
Let
\[
  A(p')=-\log(1-p'),
  \qquad
  B(p')=\log\frac{1-p'}{p'}.
\]
Then, the preceding identity can be written as
\[
  I_p(W)-I_q(W)=A(p)-A(q)+e(W)\bigl(B(p)-B(q)\bigr).
\]
Note that as $P$ is compact, both $A$ and $B$
are uniformly continuous on $P$; the singularities at $0$ and $1$ are now a
positive distance away.
Since every graphon satisfies
$0\le e(W)\le1$, we have, for all $p,q\in P$,
\[
  |I_p(W)-I_q(W)| = |A(p)-A(q)+e(W)(B(p)-B(q))|
  \le |A(p)-A(q)| + |B(p)-B(q)|.
\]
Thus, by the uniform continuity of $A$ and $B$ on $P$,
\begin{equation}\label{eq:entropy-parameter-continuity}
\sup_{W\in\W_0}|I_p(W)-I_q(W)|\to0
\qquad (|p-q|\to0,\ p,q\in P).
\end{equation}


For fixed $p_0$, lower semicontinuity of $I_{p_0}$ in the cut metric gives
$I_{p_0}(W_*)\le\liminf_n I_{p_0}(W_n)$.  Since
\eqref{eq:entropy-parameter-continuity} and $q_n\to p_0$ give
$|I_{p_0}(W_n)-I_{q_n}(W_n)|\to0$, we obtain
\begin{equation}\label{eq:entropy-moving-parameter-lsc}
  I_{p_0}(W_*)
  \le
  \liminf_{n\to\infty} I_{q_n}(W_n).
\end{equation}

On the other hand, the constant graphon $W\equiv r_n$ is feasible for
$(p_n,r_n)$.  Since $W_n$ is an optimizer,
\[
  I_{p_n}(W_n)\le J_{p_n}(r_n).
\]
Using \eqref{eq:entropy-parameter-continuity} again, now with $p=p_n$ and $q=q_n$,
and noting that $A'$ and $B'$ are bounded on $P$, we obtain
\[
  I_{q_n}(W_n)
  \le
  J_{p_n}(r_n)+O_P(|p_n-q_n|).
\]
Applying the same bound to the constant graphon $r_n$, and using
$p_n-q_n\to0$, $q_n\to p_0$, $r_n\to r_0$, and continuity of
$(p,z)\mapsto J_p(z)$, gives
\[
  J_{p_n}(r_n)=J_{q_n}(r_n)+O_P(|p_n-q_n|)\to J_{p_0}(r_0).
\]
Therefore
\begin{equation}\label{eq:boundary-cost-limsup}
  \limsup_{n\to\infty} I_{q_n}(W_n)
  \le
  J_{p_0}(r_0).
\end{equation}

Equations \eqref{eq:boundary-limit-feasibility}, \eqref{eq:entropy-moving-parameter-lsc}, and \eqref{eq:boundary-cost-limsup} show
that $W_*$ is feasible at $(p_0,r_0)$ and has $I_{p_0}$-value no larger than the
constant graphon:
\[
  I_{p_0}(W_*)\le J_{p_0}(r_0).
\]
By condition (M2) of \Cref{thm:scalar-lz-boundary},
$(r_0^d,J_{p_0}(r_0))$ lies on the convex minorant of
$x\mapsto J_{p_0}(x^{1/d})$.  Hence
Theorem~\ref{thm:lz-criterion} gives that the constant graphon $r_0$
is the unique minimizer at $(p_0,r_0)$.  Thus,  $W_*\equiv r_0$.  Since every subsequence of $(W_n)$ has a further
cut-convergent subsequence, and every such limit equals $r_0$, the whole
sequence converges to $r_0$ in the quotient cut metric.  Finally, edge
density is continuous in the cut metric, so $e(W_n)\to r_0$.
\end{proof}

\begin{proof}[Proof of \Cref{cor:scalar-reduction}]
Let $I$ and $\rho_0$ be as in \Cref{lem:fixed-density-bipodality}.
Fix $0<\rho<\rho_0$, choose $\eta>0$ as in that lemma, and take
$r\in I$, $p\in(\pc(r)-\eta,\pc(r))$, and an optimizer $W^*$ attaining
$\Phi_H(p,r)$.  By \Cref{lem:active-constraint}, $t(H,W^*)=r^m$, and by
\Cref{lem:fixed-density-bipodality},
$\eps^*=e(W^*)\in(r-\rho,r)$.  Thus, $W^*$ is admissible in the variational
problem defining $S_H(\eps^*,r^m)$, and hence
\[
  s(W^*)\le S_H(\eps^*,r^m).
\]
Hence, by \eqref{eq:relative-entropy-identity},
\[
  I_{p,r}(\eps^*)\le I_p(W^*)=\Phi_H(p,r).
\]
On the other hand, $B_{\eps^*,r^m}$ has edge density $\eps^*$ and
$H$-density $r^m$, so it is feasible for the variational problem defining
$\Phi_H(p,r)$ in \eqref{eq:graphon-variational-problem}.  Using
\eqref{eq:reduced-objective-attainment}, we get
\[
  \Phi_H(p,r)\le I_p(B_{\eps^*,r^m})=I_{p,r}(\eps^*).
\]
Thus, the preceding inequalities are all equalities, and in particular,
\[
  I_{p,r}(\eps^*)=\Phi_H(p,r),
\]
and
\[
  s(W^*)=S_H(\eps^*,r^m).
\]
To see this, suppose that the entropy inequality
$s(W^*)\le S_H(\eps^*,r^m)$ were strict.  Then the inequality
$I_{p,r}(\eps^*)\le I_p(W^*)$ would also be strict.  Therefore,
$W^*$ is an entropy maximizer under the fixed constraints
$e(W)=\eps^*$ and $t(H,W)=r^m$.  By the uniqueness part of
\Cref{thm:krrs-analytic-extension}, $W^*$ agrees with $B_{\eps^*,r^m}$ up to relabeling.

For any $\eps\in(r-\rho,r)$, the graphon $B_{\eps,r^m}$ is feasible for the
variational problem defining $\Phi_H(p,r)$ in \eqref{eq:graphon-variational-problem}, so
\[
  I_{p,r}(\eps)=I_p(B_{\eps,r^m})\ge\Phi_H(p,r)
  =I_{p,r}(\eps^*).
\]
Therefore, $\eps^*$ minimizes $I_{p,r}$ on $(r-\rho,r)$, and
\eqref{eq:edge-density-minimization} follows.  The same inequality gives the converse
statement: if $\eps$ minimizes $I_{p,r}$ on $(r-\rho,r)$, then
$I_p(B_{\eps,r^m})=I_{p,r}(\eps)=I_{p,r}(\eps^*)
=\Phi_H(p,r)$, and $B_{\eps,r^m}$ is feasible, so it is an optimizer
attaining $\Phi_H(p,r)$.
Finally, distinct $\eps$ give graphons of distinct edge density, hence
inequivalent under relabeling, which gives the asserted bijection.
\end{proof}

\subsection{Parameter asymptotics near nonexceptional endpoints}
\label{app:bipodal-parameter-expansions}

\begin{remark}[First-order asymptotics of the bipodal parameters]
\label{rmk:bipodal-parameter-expansions}
\Cref{thm:nonexceptional-endpoint} gives the qualitative limit
\(c(p,r)\to0\).  The purpose of this remark is to quantify that limit and the
corresponding change in the block densities.  Work in the neighborhood $U$
of that theorem, shrinking it if necessary, and let \(p\uparrow\pc(r)\)
from the symmetry-breaking side.  Relabel the podes as in \Cref{thm:nonexceptional-endpoint}, so
that the first pode has size \(c=c(p,r)\), and set
\[
  \delta_*:=r-e(W_{p,r}),
  \qquad
  z:=\zeta_d(r).
\]
Denote by \(q_{11}^0(r)\) the value at
\((\eps,\vartheta)=(r,0)\) of the analytically extended KRR--S parameter
\(q_{11}\) from \Cref{thm:krrs-analytic-extension}.
All limits below refer to this regime, and the $O_{H,U}$-bounds are uniform
for $(p,r)\in U$ with $p<\pc(r)$.

For \(z\ne r\), define
\begin{equation}\label{eq:edge-deficit-coefficient}
  D_d(r,z):=\frac{z^d-r^d}{d r^{d-1}}-(z-r)>0,
\end{equation}
where positivity follows from the strict convexity of \(z\mapsto z^d\).
Since $\zeta_d(r_0)\ne r_0$, continuity lets us shrink $U$ so that
$D_d(r,\zeta_d(r))$ is uniformly bounded away from zero.
Then
\[
  \delta_*=2D_d(r,z)c+O_{H,U}(\delta_*^2).
\]
Thus, \(2D_d(r,z)\) is the first-order edge-density deficit per unit of
smaller-pode mass, after \(q_{22}\) is adjusted to preserve the
\(H\)-density.  Combining this relation with
\(\delta_*=\lambda(p,r)/A_H(r)+O_{H,U}(\lambda(p,r)^2)\) from
\Cref{thm:nonexceptional-endpoint} gives
\begin{equation}\label{eq:optimizer-block-size}
  c(p,r)=
  \frac{\lambda(p,r)}
  {2D_d(r,\zeta_d(r))A_H(r)}
  +O_{H,U}(\lambda(p,r)^2),
\end{equation}
\begin{equation}\label{eq:optimizer-block-density}
  q_{22}(p,r)
  =
  r-
  \frac{\zeta_d(r)^d-r^d}
       {d r^{d-1}D_d(r,\zeta_d(r))A_H(r)}
  \lambda(p,r)
  +O_{H,U}(\lambda(p,r)^2),
\end{equation}
and
\[
  q_{12}(p,r)=\zeta_d(r)+O_{H,U}(\lambda(p,r)),
  \qquad
  q_{11}(p,r)=q_{11}^0(r)+O_{H,U}(\lambda(p,r)).
\]
In particular, both the smaller-pode size and the displacement of the
large-pode density from \(r\) are of order \(\lambda(p,r)\).
The other two block densities satisfy
\(q_{12}\to\zeta_d(r)\ne r\) and
\(q_{11}\to q_{11}^0(r)\).  These relations quantify the vanishing-pode
mechanism described in the introduction.

Although they arise from different formulations, the KRR--S limiting cross
density and the second contact density of the Lubetzky--Zhao supporting line
coincide: \(\zeta_d(r)=\sm(r)\) for every \(r\ne r_*\).
Indeed, put \(\beta:=J_{\pc(r)}'(r)/(d r^{d-1})\).
Since \(J_{\pc(r)}+2S_0\) is affine,
\eqref{eq:supporting-cost-gap} and \eqref{eq:edge-deficit-coefficient} give
\[
  g_r(z)=-d r^{d-1}D_d(r,z)\bigl\{\psi_d(r,z)+\beta\bigr\}
  \qquad(z\in(0,1)\setminus\{r\}).
\]
The positivity of \(D_d(r,z)\) and \(g_r\ge0\) imply
\(\psi_d(r,z)\le-\beta\), also at \(z=r\) by continuity.
Equality holds at \(z=\sm(r)\) by \Cref{thm:scalar-lz-boundary}, so
the uniqueness of the maximizer in \Cref{thm:krrs-cross-density} proves the
identity.  Thus, the comparison graphons in \Cref{sec:nonexceptional-quadratic-growth} use
exactly the limiting cross density of the optimizer.

For completeness, we derive the two explicit expansions above.  Suppress the
arguments \((p,r)\) and evaluate the analytic KRR--S parameter map at
\[
  (\eps,\vartheta)
  =\bigl(r-\delta_*,\,r^m-(r-\delta_*)^m\bigr).
\]
After shrinking $U$, the segments from $(r,0)$ to these parameter pairs
lie in a compact subset of the KRR--S analytic domain, as in
\Cref{sec:nonexceptional-proof}.  Uniform Taylor estimates, the boundary
values in \Cref{thm:krrs-analytic-extension}, and
\(\vartheta=O_{H,U}(\delta_*)\) therefore give
\[
  c=O_{H,U}(\delta_*),\qquad
  q_{22}-r=O_{H,U}(\delta_*),\qquad
  q_{12}-z=O_{H,U}(\delta_*),\qquad
  q_{11}-q_{11}^0(r)=O_{H,U}(\delta_*).
\]
The exact edge-density identity
\[
  e(W_{p,r})
  =q_{22}+2c(q_{12}-q_{22})
   +c^2(q_{11}-2q_{12}+q_{22})
\]
therefore yields
\begin{equation}\label{eq:bipodal-edge-linearization}
  -\delta_*=(q_{22}-r)+2(z-r)c+O_{H,U}(\delta_*^2).
\end{equation}

For the \(H\)-density, separate the homomorphisms according to the number of
vertices mapped to the smaller pode.  Since \(H\) is \(d\)-regular and has
\(2m/d\) vertices, the contributions with zero and one such vertex give
\begin{align*}
  t(H,W_{p,r})
  &=(1-c)^{2m/d}q_{22}^m
    +\frac{2m}{d}c(1-c)^{2m/d-1}
      q_{12}^d q_{22}^{m-d}+O_{H,U}(c^2)\\
  &=r^m+m r^{m-1}(q_{22}-r)
    +\frac{2m}{d}r^{m-d}(z^d-r^d)c
    +O_{H,U}(\delta_*^2).
\end{align*}
Using \(t(H,W_{p,r})=r^m\), we obtain
\begin{equation}\label{eq:bipodal-density-linearization}
  0=m r^{m-1}(q_{22}-r)
    +\frac{2m}{d}r^{m-d}(z^d-r^d)c
    +O_{H,U}(\delta_*^2).
\end{equation}
Eliminating \(q_{22}-r\) between
\eqref{eq:bipodal-edge-linearization} and
\eqref{eq:bipodal-density-linearization} gives
\[
  \delta_*=2D_d(r,z)c+O_{H,U}(\delta_*^2),
  \qquad
  c=\frac{\delta_*}{2D_d(r,z)}+O_{H,U}(\delta_*^2),
\]
which proves \eqref{eq:optimizer-block-size}.  Finally,
\eqref{eq:bipodal-edge-linearization} gives
\[
  q_{22}
  =r-\frac{z^d-r^d}{d r^{d-1}D_d(r,z)}\delta_*
   +O_{H,U}(\delta_*^2),
\]
and \eqref{eq:optimizer-block-density} follows.

When $d=2$, a direct computation using $S_0(1-r)=S_0(r)$ and
$S_0'(1-r)=-S_0'(r)$ gives $\zeta_2(r)=1-r$, which yields the
triangle specialization of \eqref{eq:optimizer-block-size} used in the
introduction.
\end{remark}
